% INTERNAL WORKING COPY: do not submit until all bold placeholders are closed and the manuscript Diff is applied and checked.
% Source-alignment note: existing manuscript-edit claims outside the neutrino response (including the Introduction wording) refer to the separate Route B revision copy documented with PR #13. The neutrino replies follow the transcript-aligned response v2 and propose a short manuscript clarification without claiming that it has already been applied. The canonical original main-v6.tex has not been overwritten in this repository. Do not treat the response as accepted until the final manuscript source baseline is selected, synchronized, and all remaining placeholders and reference dependencies are resolved.
\documentclass[12pt,a4paper]{article}

% Mathematics
\usepackage{amsmath,amssymb}
\usepackage{bbm}
\usepackage{slashed}

% Graphics
\usepackage{graphicx}
\usepackage{epsfig}
\usepackage{epstopdf}

% Tables and page layout
\usepackage{multirow}
\usepackage{placeins}
\setlength{\parskip}{10pt}

% Text formatting and URLs
\usepackage[normalem]{ulem}
\usepackage{url}

% Colors and cross-document citations
\usepackage[dvipsnames]{xcolor}
\usepackage{xcite}
\externalcitedocument{main-v6}

% Hyperlinks (load near the end)
\usepackage[
    colorlinks=true,
    citecolor=purple,
    linkcolor=blue
]{hyperref}

% Custom commands
\newcommand{\GJH}[1]{\textcolor{blue}{GJH: #1}}
\newcommand{\WXP}[1]{\textcolor{magenta}{#1}}

\begin{document}

%\title{Referee Report on DY13253}
%\author{Rundong Fang, Jinhui Guo, Ming Li, Jia Liu, Xiao-Ping Wang, Yiheng Xiong}
%\date{2026.19.22}
%\maketitle

\noindent \textbf{Replies to Referee}
\medskip

\subsection*{Replies to Referee}

\noindent Dear Editor,

This paper considers an anomaly-free, renormalizable model of an additional Abelian
gauge boson \(A'\) with chiral couplings to muons and taus. In this model, vector-like
leptons and a scalar field \(S\) are introduced to construct the ultraviolet completion
of the axial \(U_A(1)_{L_{\mu-\tau}}\) gauge theory. The phenomenology of the \(A'\) in
the axial limit is studied, and the relevant experimental constraints are discussed.
Additionally, possible collider signatures of the model are investigated.

The referee summarizes our work as follows:

\begin{quote}
\itshape
The authors construct an anomaly-free, renormalizable model in which an additional Abelian gauge field \(A'\) has chiral couplings to muons and taus. This model realizes a limit where the \(A'\) coupling to the \(\mu-\tau\) current is axial. The second goal of the paper is to examine the phenomenology of this model in the axial limit.

The paper presents results that are useful for the model building community. However, I have some significant reservations about key points that either should be included or otherwise justified in being excluded from the manuscript. These are explained below. Revisions for these major points (even if only to clarify why they are not issues) are required before the paper is at the level of a PRD manuscript.
\end{quote}

We sincerely thank the referee for the positive assessment of the model-building construction and for identifying issues that require further clarification or quantitative analysis. We respond to each point below.

\medskip
\noindent {\bf Major Points 1)}. \begin{quote}\itshape
Phenomenology of the \(S\) scalar. The second half of the paper reviews a set of experimental bounds on only an effective theory of the \(A'\) axially coupled to the \(L_{\mu-\tau}\) current. This leaves a major gap between these two parts of the paper: the authors show that their UV theory requires a scalar \(S\) that cannot be arbitrarily decoupled from the low-energy dynamics, e.g. from Eq. (34) which states that \(m_S^2 < 4\pi\left(m_{A'}/(2g_X)\right)^2\).\end{quote}

\noindent {\bf Reply.}
We thank the referee for highlighting the need to clarify the role of the scalar \(S\) in the UV completion and its phenomenological implications. We agree that \(S\) is an intrinsic component of the symmetry-breaking sector, and that its phenomenological relevance must be assessed under clearly stated assumptions and in light of the applicable theoretical and experimental constraints.

For the small-portal reference benchmark used in the Fig.~2 analysis, the scalar mass relations are
\[
m_{A'}=2g_Xv_s,\qquad m_s^2=2\lambda_Sv_s^2,
\qquad
\lambda_S=\frac{2g_X^2m_s^2}{m_{A'}^2}.
\]
Imposing the perturbativity criterion \(\lambda_S\leq\lambda_S^{\max}\) therefore gives
\[
m_s\leq\sqrt{2\lambda_S^{\max}}\,v_s
=\sqrt{\frac{\lambda_S^{\max}}{2}}\frac{m_{A'}}{g_X}.
\]

For the adopted mass convention, the perturbativity condition is
\[
\lambda_S\leq\lambda_S^{\max}
\quad\Longrightarrow\quad
m_s^2\leq 2\lambda_S^{\max}v_s^2.
\]
To make the coefficient in the referee's quoted Eq.~(34) unambiguous, \(m_s^2<4\pi v_s^2\) corresponds to \(\lambda_S^{\max}=2\pi\) in this potential convention, whereas the \(\lambda_S^{\max}=4\pi\) benchmark used in our numerical example gives \(m_s^2\leq8\pi v_s^2\), or \(m_s\leq\sqrt{8\pi}\,v_s\). We will check the coefficient against the exact submitted version and state the chosen perturbativity threshold explicitly, using a consistent convention in Eq.~(34), the numerical example, and the figure caption. For \(m_s=10\,\mathrm{GeV}\), this latter benchmark gives \(v_s\gtrsim2.0\,\mathrm{GeV}\). In particular, \(m_s\leq v_s\) is not a general consequence of \(\lambda_S\leq4\pi\); it would require the stronger condition \(\lambda_S\leq1/2\). Conversely, within the plotted range \(g_X\leq0.1\), the vector mass relation gives \(m_{A'}=2g_Xv_s\leq0.2v_s\). These are conditional theoretical relations, not experimental exclusions, and do not by themselves establish an experimentally allowed scalar-mass interval.

For clarity, we reproduce here the selected, non-exhaustive low-mass scalar constraints listed in Table~II of the manuscript. They constrain the scalar coupling \(y_{s\mu\mu}\), not the gauge-boson coupling \(g_X\) directly:
\begin{center}
\small
\begin{tabular}{p{0.28\linewidth}p{0.39\linewidth}p{0.20\linewidth}}
\hline
Scalar-mass range & Published constraint & Approximate upper bound on \(|y_{s\mu\mu}|\)\\
\hline
\(m_s<10^{-10}\,\mathrm{eV}\) & Neutron-star binaries~\cite{Dror:2019uea,Liu:2025zuz} & \(\lesssim10^{-20}\)\\
\(m_s\lesssim10^{-2}\,\mathrm{eV}\) & Fifth-force tests~\cite{Hoskins:1985tn,Kapner:2006si,Blakemore:2021zna} & \(\lesssim10^{-16}\)\\
\(m_s\lesssim100\,\mathrm{keV}\) & Neutron-star cooling~\cite{Fiorillo:2026wso} & \(\lesssim1.2\times10^{-12}\)\\
\(m_s\lesssim\mathcal{O}(10)\,\mathrm{MeV}\) & SN~1987A cooling~\cite{Caputo:2021rux,Balazs:2022tjl} & \(\lesssim4.2\times10^{-9}\)\\
\(100\,\mathrm{keV}\lesssim m_s\lesssim\mathcal{O}(10)\,\mathrm{MeV}\) & SN~1987A and diffuse \(\gamma\)-rays~\cite{Caputo:2021rux,Balazs:2022tjl} & \(\lesssim5\times10^{-11}\)\\
\hline
\end{tabular}
\end{center}
These are approximate, source-dependent bounds on \(y_{s\mu\mu}\). They inherit the production, trapping, and loop-matching assumptions of the cited analyses; they are not an exhaustive global constraint set and do not establish the strongest bound at every scalar mass. They are not direct, model-independent bounds on \(g_X\). The low-mass scalar bounds reproduced here are distinct from the collider and \((g-2)_\mu\) constraints discussed for the \(A'\) effective current; the latter do not, by themselves, establish a scalar-specific exclusion. We keep these constraint categories separate rather than presenting Fig.~2 as a joint scalar--\(A'\) scan. Only under the mass-proportional benchmark \(y_{s\mu\mu}=m_\mu/v_s\), together with \(m_{A'}=2g_Xv_s\), can one translate a bound \(|y_{s\mu\mu}|<y_{\max}\) into the conditional relation
\[
g_X<\frac{y_{\max}m_{A'}}{2m_\mu}.
\]
This is the same conditional conversion used in Eq.~(31) of the original manuscript. It is not a model-independent bound on \(g_X\); it applies only when the scalar--muon Yukawa normalization and benchmark assumptions stated above are satisfied.

The tabulated limits are a summary of selected low-mass constraints on the scalar--muon coupling, not a complete scalar-sector exclusion map. They inherit the production, trapping, and matching assumptions of the cited analyses. Collider constraints additionally depend on the model-specific production rates, portal-induced mixing, branching fractions (including \(s\to A'A'\) when kinematically open), widths, and lifetime. Since a full model-specific scalar recast has not been completed, we do not classify the full scalar-mass domain as allowed or excluded on the basis of the mass relation or published search coverage alone.

\medskip
\noindent {\bf Major Points 1.1)}. \begin{quote}\itshape
It is thus an odd choice to exclude the scalar from the phenomenology of the effective theory: if the \(S\) is unavoidable at phenomenologically relevant scales, shouldn't one of the main points of the paper be that the only known UV completion of an axial \(L_{\mu-\tau}\) gauge boson has a prediction for a scalar particle in a mass range that can be probed by colliders?\end{quote}

\noindent {\bf Reply.}
We agree that \(S\) is a genuine prediction of the UV construction. The purpose of this paper, however, is to study the axial-current phenomenology of \(A'\), not to provide a comprehensive collider analysis of every singlet-like scalar realization. This scope choice does not mean that \(S\) is absent or can be arbitrarily decoupled; its role in symmetry breaking, its mass relation, and the applicable constraints are addressed in this response.

The generic phenomenology of singlet-like scalars mixed with the Standard Model Higgs—including Higgs signal-strength constraints, direct searches, and the dependence on scalar masses and mixing—is well established; see Robens and Stefaniak (\href{https://arxiv.org/abs/1501.02234}{arXiv:1501.02234}) and Robens, Stefaniak, and Wittbrodt (\href{https://arxiv.org/abs/1908.08554}{arXiv:1908.08554}). The GeV-scale scalar-portal phenomenology and its production channels have also been studied systematically by Boiarska et al. (\href{https://arxiv.org/abs/1904.10447}{arXiv:1904.10447}). We cite these works to place our scope choice in context, rather than repeat a comprehensive generic scalar-portal analysis in a paper whose principal target is the axial-vector \(A'\) phenomenology. The published numerical limits cannot, however, be transferred directly to our UV model without matching the couplings, branching fractions, widths, lifetimes, and production rates, including \(s\to A'A'\) when kinematically open.

Dedicated searches for leptophilic scalars have also been performed at BaBar, Belle, and Belle II \cite{BaBar:2020jma,Belle:2022gbl,Belle-II:2024wtd,Belle-II:2023ydz}. In particular, the Belle II four-muon resonance search (arXiv:2403.02841) covers approximately \(0.212\text{--}9.0\,\mathrm{GeV}\) and includes a muonphilic spin-0 interpretation, while the \(e^+e^-\to\mu^+\mu^-\tau^+\tau^-\) search (arXiv:2306.12294) covers approximately \(3.6\text{--}10\,\mathrm{GeV}\) and includes a Higgs-like leptophilic spin-0 interpretation. The associated-production search and the four-muon search probe different production channels and final states, so they cannot be combined into a universal scalar limit. For heavier masses, prospective sensitivity at high-energy colliders, including a future muon collider, has been studied for leptophilic-scalar benchmarks with different Yukawa and CP assumptions \cite{Chen:2021leptophilic}. This provides context for possible future searches, not a reach projection for the radial scalar in our UV model. Mapping existing searches to our UV scalar requires matching the coupling normalization and evaluating production, branching fractions, widths, lifetimes, and acceptance at the same UV parameter point, including \(s\to A'A'\) when kinematically open. This model-specific recast has not been completed; published mass coverage alone therefore does not establish an exclusion or discovery reach, and the scalar regions not covered by an applicable model-specific analysis remain unclassified.

\medskip
\noindent {\bf Major Points 1.2)}. \begin{quote}\itshape
What is the range of phenomenologically allowed \(S\) mass including perturbativity limits on \(\lambda\) and a motivated range of \(m_{A'}\) and \(g_X\). For example: perhaps the authors imagine a range of \(m_{A'}\) below some benchmark mass where certain experimental constraints on \(g_X\) become very strong. Then with that set of constraints (which should be stated clearly, see minor point \#1 below), one may give a range of \(m_S\) and comment on the experimental constraints.\end{quote}

\noindent {\bf Reply.}
The scan window in Fig.~2 is
\[
10^{-3}\,\mathrm{GeV}\leq m_{A'}\leq10^2\,\mathrm{GeV},
\qquad
10^{-6}\leq g_X\leq10^{-1}.
\]
These intervals define the plotted domain; they are not universal theoretical boundaries and do not imply that every point inside the window is viable. In the minimal one-VEV benchmark, \(v_s=m_{A'}/(2g_X)\) is derived at each scan point, and the scalar self-coupling follows from
\[
\lambda_S=\frac{m_s^2}{2v_s^2}
=\frac{2g_X^2m_s^2}{m_{A'}^2}.
\]
The perturbativity criterion \(\lambda_S\leq\lambda_S^{\max}\) therefore constrains \(m_s/v_s\), equivalently the scalar self-coupling, for the chosen benchmark. Since \(v_s=m_{A'}/(2g_X)\) is derived point by point, this relation does not by itself give a unique scalar-mass interval or an independent universal exclusion on \(m_{A'}\) or \(g_X\); the phenomenological classification also depends on the applicable constraints and the specified scalar couplings and decay pattern.

The mass projection need not be a single continuous interval. The final classification must distinguish regions excluded by the constraints actually applied, regions not excluded by those constraints, and regions not yet assessed. Thus, at each selected \((m_{A'},g_X)\), perturbativity supplies a conditional theoretical upper bound on \(m_s\), while the phenomenologically allowed set is determined only after applying the constraints that are valid for the specified scalar couplings, portal mixing, and decay pattern. The low-mass coupling limits reproduced above do not determine a complete allowed mass interval, and a full collider/portal recast has not been completed. We therefore do not infer that the full scan window is viable or present a single universal interval for \(m_s\); instead, the plotting window, theoretical consistency condition, regions excluded by constraints actually applied, and regions not yet assessed are kept distinct.

Satisfying the perturbativity criterion alone is not taken as evidence of experimental viability.

\medskip
\noindent {\bf Major Points 1.3)}. \begin{quote}\itshape
Above (34): the authors ``assume negligible Higgs-portal mixing.'' Is this interaction not generated at loop level? It seems that \(S\)-\(h\) mixing cannot be avoided due to \(\mu\) and \(\tau\) loops. In turn, this seems like it would lead to strong experimental bounds on the UV completion. In fact, once there is mixing, is there not a resonance where the mixing becomes large as \(m_S\to m_h\)? It would be useful to quantify the size of this mixing over the allowed range of \(m_S\). Existing Higgs measurements may set the strongest bounds on the UV completion. If these bounds are obviously negligible (though again, I expect this to be strongly constrained as \(m_S\to m_h\)), then the authors should present this as the justification for ignoring \(m_S\) in the low-energy theory. (Though I am still skeptical that the \(S\) can be ignored since it cannot formally decouple.)\end{quote}

\noindent {\bf Reply.}
We agree that a vanishing Higgs-portal coupling is not enforced by the gauge symmetry: the operator \((H^\dagger H)(S^\dagger S)\) is allowed, so the portal boundary condition must be stated at a matching scale and in a specified renormalization convention. The negligible-portal assumption is therefore used only as a conditional phenomenological benchmark, not as a symmetry-protected statement or a claim that the portal vanishes at every scale.

The general phenomenology of Higgs--singlet mixing, including its impact on the \(125\) GeV Higgs signal strengths and searches for additional scalar states, has already been studied in detail; representative analyses are cited above. This is one reason we do not reproduce a full generic singlet-Higgs analysis in a paper whose main objective is the axial-vector \(A'\) phenomenology. These studies do not, however, determine the physical mixing in our particular UV model, and we do not use them as a substitute for the model-specific loop check.

The zero-momentum effective-potential analysis for the vector-like-lepton (VLL) sector gives, under its stated potential convention,
\[
\lambda_{HS}^{\mathrm{EFT}}(\mu_R)
=\lambda_{HS}^{\mathrm{UV}}(\mu_R)
-\sum_{\ell=\mu,\tau}
\frac{|y_\ell|^2|\lambda_\ell|^2}{8\pi^2}
+\cdots,
\]
where \(y_\ell\) and \(\lambda_\ell\) are the Yukawa couplings in the VLL interaction. This is a conditional zero-momentum threshold relation; it does not determine the finite-momentum off-diagonal self-energy or a physical scalar mixing angle.

To make the portal assumption explicit, we have evaluated the leading VLL threshold at the matching scale \(\mu_M=M_E\),
\[
\Delta\lambda_{HS}^{E}
=-\frac{1}{8\pi^2}\sum_{\ell=\mu,\tau}|y_\ell|^2|\lambda_\ell|^2,
\qquad
\lambda_{HS}^{\mathrm{EFT}}(\mu_M)
=\lambda_{HS}^{\mathrm{UV}}(\mu_M)+\Delta\lambda_{HS}^{E}.
\]
We use \(\lambda_{HS}^{\mathrm{EFT}}(M_E)=0\) as an explicit small-portal reference condition, which requires a nonzero UV portal \(\lambda_{HS}^{\mathrm{UV}}=-\Delta\lambda_{HS}^{E}\). This condition is a benchmark choice, not a symmetry prediction. For illustration, taking \(M_E=20\,\mathrm{TeV}\), \(\theta_{\mu L}=\theta_{\tau L}=10^{-2}\), and the leading seesaw relations gives \(\Delta\lambda_{HS}^{E}\simeq-0.2108\) at \((m_{A'},g_X,v_s)=(10\,\mathrm{GeV},0.0705,70.9\,\mathrm{GeV})\), and \(\Delta\lambda_{HS}^{E}\simeq-0.01067\) at \((63\,\mathrm{GeV},0.1,315\,\mathrm{GeV})\). For these illustrative points only, we additionally impose \(|\lambda_\ell|\leq\sqrt{4\pi}\); this is a theory-consistency benchmark rather than a new exclusion contour in Fig.~2. The comparison demonstrates that setting the UV portal to zero is not equivalent to setting the matched EFT portal to zero, and that the latter must be stated as an input if the small-portal benchmark is retained.

In the present zero-tree-level-mixing setup, the singlet-like field \(s\) has no tree-level \(sWW\) or \(s\bar t t\) vertex; accordingly, direct one-loop \(W\)- or top-quark diagrams do not contribute to \(\Pi_{hs}\) at this order. The leading one-loop off-diagonal contribution instead comes from the charged-lepton/vector-like-lepton sector. We have evaluated this fermion contribution using the full \(2\times2\) charged-lepton/VLL mass-eigenstate matrices, including light-light, light-heavy, heavy-light, and heavy-heavy internal states, and have separated the light-lepton EFT contribution from the hard momentum-dependent term. We use the explicit reference condition \(\lambda_{HS}^{\mathrm{EFT}}(M_E)=0\) stated above; the result is therefore conditional on that matching choice and is not a prediction independent of the renormalized portal boundary.

For illustration, with \(M_E=20\,\mathrm{TeV}\), \(\theta_{\mu L}=\theta_{\tau L}=10^{-2}\), and the same leading-seesaw benchmark inputs used for the threshold check, the matched off-diagonal self-energy at \(p^2=m_h^2\) is
\[
\Pi_{hs}^{\mathrm{matched}}(m_h^2)
\simeq(-0.4133-0.1132i)\,\mathrm{GeV}^2
\]
for \((m_{A'},g_X,v_s)=(10\,\mathrm{GeV},0.0705,70.9\,\mathrm{GeV})\), and
\[
\Pi_{hs}^{\mathrm{matched}}(m_h^2)
\simeq(-0.09300-0.02547i)\,\mathrm{GeV}^2
\]
for \((m_{A'},g_X,v_s)=(63\,\mathrm{GeV},0.1,315\,\mathrm{GeV})\). The exact mass-eigenstate hard momentum difference reproduces the leading finite-derivative matching result to \(6.1\times10^{-4}\) relative in both benchmarks, while the zero-momentum threshold comparison agrees to \(1.8\times10^{-4}\) relative under the leading-seesaw input convention.

The finite-momentum result does not imply uniformly negligible mixing. At \(m_s=m_h=125.25\,\mathrm{GeV}\), the width-regulated diagnostic \(|\Pi_{hs}/(a-b)|\), with \(a=m_h^2-i m_h\Gamma_h\) and \(b=m_s^2-i m_s\Gamma_s\), is approximately \(9.1\times10^{-4}\) when \(s\to A'A'\) is open and \(\Gamma_s\simeq3.75\,\mathrm{GeV}\). For the benchmark in which \(s\to A'A'\) is closed and \(\Gamma_s\simeq0.159\,\mathrm{MeV}\), the same diagnostic reaches about \(0.197\) at exact degeneracy; the corresponding \(2|\Pi_{hs}/(a-b)|\) diagnostic is about \(0.394\). Thus the width difference regulates the open-channel case, but a narrow near-degenerate region in the closed-channel case can have non-negligible mixing.

On a \(1\,\mathrm{MeV}\)-grid diagnostic scan over the illustrative interval \(m_s\in[10,150]\,\mathrm{GeV}\), the criterion \(2|\Pi_{hs}/(a-b)|\geq0.1\) is never reached in the open-channel benchmark. In the closed-channel benchmark, the diagnostic crosses this threshold only for \(m_s\in[125.24255,125.25745]\,\mathrm{GeV}\), corresponding to a half-width of \(7.446\,\mathrm{MeV}\) around \(m_h=125.25\,\mathrm{GeV}\). This scan holds the matched \(\Pi_{hs}(m_h^2)\) fixed at its benchmark value while varying \(m_s\) and the tree-level singlet width; it is therefore a diagnostic approximation, not a self-consistent two-pole solution or an experimental exclusion interval.

These are benchmark-level one-loop fermion results with the stated matching condition, not a global scan of the UV model. The narrow-width near-degenerate region therefore remains \textbf{not assessed}: a complete pole-scheme treatment and model-specific two-pole Higgs lineshape/signal-strength and direct-search recasts are still needed before assigning allowed/excluded status. We do not use these benchmark results to justify neglecting \(S\) in the full UV model or to classify the full UV parameter space. The scalar \(S\) remains in the UV construction, and the \(A'\) contours remain conditional constraints on the specified small-portal, pure-axial effective-current benchmark.

\medskip
\noindent {\bf Major Points 2)}. \begin{quote}\itshape
Motivating the range of effective theory parameters. To further aid in the comparison between the UV theory and the two-parameter low-energy theory, it would be useful to explicitly state in one place (perhaps in a table or a figure like Fig. 2) the range of the effective model parameters that are feasible. The main phenomenological results are in Fig. 2, but it was not clear why this range of couplings and \(A'\) masses were chosen: are values outside this range not allowed? What, if anything, does the perturbativity of the UV theory impose on the effective parameters?\end{quote}

\noindent {\bf Reply.}
We thank the referee for asking us to motivate the parameter window and to distinguish the low-energy effective-current analysis from the UV theory. The ranges displayed in Fig.~2 are
\[
10^{-3}\,\mathrm{GeV}\leq m_{A'}\leq10^2\,\mathrm{GeV},
\qquad
10^{-6}\leq g_X\leq10^{-1}.
\]
They define the plotting window adopted for this comparison, not universal theoretical or experimental boundaries. Values outside this window are not excluded by the choice of axes alone, and points inside it are not automatically viable in the complete UV model.

The window is intended to display the broad mass interval relevant to the low-energy precision constraints and the terrestrial searches discussed in this work. The meeting discussion identified the MeV-to-\(100\,\mathrm{GeV}\) interval as a practical range for this comparison: the low-mass end includes the region where low-energy precision constraints, including the stated \((g-2)_\mu\) comparison, are relevant, while the upper end reaches the terrestrial search coverage considered in the paper. The stellar and white-dwarf bounds on the scalar listed separately in Table~II constrain \(y_{s\mu\mu}\) under their own assumptions; they should not be read as direct constraints on the \(A'\) effective-current coupling \(g_X\). We therefore keep scalar-sector bounds and \(A'\) effective-current bounds distinct. The displayed window could be extended in either direction for a different study, but its endpoints are chosen for the scope and coverage of this analysis, not because all lower or higher masses are forbidden. We also do not claim that any single constraint is strongest throughout the full plane.

For clarity, the scan variables and their UV interpretation are summarized here:
\begin{center}
\small
\begin{tabular}{p{0.19\linewidth}p{0.34\linewidth}p{0.39\linewidth}}
\hline
Quantity & Range or relation & Interpretation\\
\hline
\(m_{A'}\) & \(10^{-3}\)--\(10^2\,\mathrm{GeV}\) & Chosen plotting window; not a universal mass bound.\\
\(g_X\) & \(10^{-6}\)--\(10^{-1}\) & Chosen plotting window; the upper edge is well below the usual perturbative criterion \(g_X\lesssim\sqrt{4\pi}\). The lower edge is not a universal lower bound.\\
\(v_s\) & \(v_s=m_{A'}/(2g_X)\) & Derived from each selected \((m_{A'},g_X)\) point, not an independent scan variable.\\
\(\lambda_S\) & \(\lambda_S=m_s^2/(2v_s^2)=2g_X^2m_s^2/m_{A'}^2\) & Scalar self-coupling in the stated small-portal benchmark; its perturbativity condition must be checked separately.\\
\(m_s\) benchmark & \(m_s\geq10\,\mathrm{GeV}\) for the scalar-suppressed reference case & Conditional benchmark choice, not a claim that \(s\) is automatically decoupled or experimentally allowed.\\
\hline
\end{tabular}
\end{center}

In the scalar-mass convention used in the manuscript, imposing \(\lambda_S\leq\lambda_S^{\max}\) gives
\[
m_s\leq\sqrt{2\lambda_S^{\max}}\,v_s.
\]
For the illustrative choice \(\lambda_S^{\max}=4\pi\), this becomes \(m_s\leq\sqrt{8\pi}\,v_s\). Consequently, the reference benchmark \(m_s\geq10\,\mathrm{GeV}\) requires
\[
v_s\geq\frac{10\,\mathrm{GeV}}{\sqrt{8\pi}}
\simeq1.995\,\mathrm{GeV},
\qquad
m_{A'}\gtrsim(3.99\,\mathrm{GeV})\,g_X.
\]
This is a conditional theoretical consistency requirement for that scalar benchmark. A point that fails it does not support the chosen \(m_s\geq10\,\mathrm{GeV}\), \(\lambda_S\leq4\pi\) benchmark; it is not thereby excluded from every possible realization of the UV model. Within the displayed range \(g_X\leq0.1\), the gauge-coupling perturbativity criterion is not the limiting condition. The scalar self-coupling condition instead imposes a point-dependent check once \(m_s\) is specified.

The effective-current contours in Fig.~2 are conditional constraints in the stated pure-axial benchmark. The final UV interpretation also depends on the portal boundary condition, finite vector-like-lepton effects, kinetic mixing, EFT validity and truncation, and the applicability of the relevant experimental constraints at the same UV parameter point. These checks do not yet provide a global classification of the full UV parameter space. We therefore distinguish the plotting window, theoretical consistency conditions, regions excluded by constraints actually applied, regions not excluded by those constraints, and regions that remain unassessed. Satisfying perturbativity alone is not evidence of experimental viability.

\medskip
\noindent {\bf Major Points 3)}.

\begin{quote}\itshape
Are neutrino oscillations possible? The Standard Model is defined with only left-handed neutrinos because of the possibility that neutrinos are their own antiparticle. In such a case neutrino masses are generated through, e.g. through the Weinberg operator.
\end{quote}

\medskip
\noindent {\bf Reply to Point 3.1}

\begin{quote}\itshape
In this model, \(U_A(1)_{L_{\mu-\tau}}\) invariance allows Weinberg operators for each generation. However, it appears that both the Yukawa matrices and neutrino masses must be aligned and therefore a non-trivial PMNS matrix does not appear to be possible.
\end{quote}

\noindent {\bf Reply.}

We thank the referee for this comment. The focus of the present work is the construction and phenomenology of the axial \(U_A(1)_{L_\mu-L_\tau}\) interaction. We have not attempted to provide a simultaneous explanation of the neutrino mass spectrum and the full PMNS matrix. Nevertheless, we agree that the neutrino-mass terms allowed by the present field content, and the limitations of the resulting mixing structure, should be stated explicitly.

With the charge assignment in Table~I, namely
\[
q(L_e,L_\mu,L_\tau)=(0,1,-1),\qquad q(H)=0,\qquad q(S)=2,
\]
where \(q=Q/Q_\ell\), Majorana neutrino masses can be introduced through symmetry-allowed effective operators. Writing \(\mathcal O_{ij}=(L_i\!\cdot H)(L_j\!\cdot H)\) with the Lorentz contraction understood, the dimension-five operators allow the \(ee\) and \(\mu\tau\) entries. The \(\mu\mu\) and \(\tau\tau\) entries require the dimension-six operators \(S^\dagger\mathcal O_{\mu\mu}/\Lambda_\nu^2\) and \(S\mathcal O_{\tau\tau}/\Lambda_\nu^2\), respectively. After symmetry breaking, the resulting mass matrix has the form
\[
M_\nu=
\begin{pmatrix}
a&0&0\\
0&b&c\\
0&c&d
\end{pmatrix},
\]
where \(a,c\) arise from dimension-five operators and \(b,d\) can arise from the operators containing \(S^\dagger,S\).

This structure permits nonzero neutrino masses and mixing within the \(\mu\)--\(\tau\) sector. However, it does not generate the electron--muon and electron--tau mixing required for the full observed PMNS matrix. The missing entries cannot be generated by further insertions of \(S\) alone: its charge-two vacuum expectation value preserves a residual \(Z_2\) symmetry separating the electron sector from the muon and tau sectors. The charged-lepton mass matrix respects the same separation. Thus, complete alignment and an identity PMNS matrix are not required by the symmetry, but the electron-flavour mixing remains absent. The present field content therefore allows a restricted neutrino-mass structure, without reproducing the complete observed mixing pattern.

\medskip
\noindent {\bf Reply to Point 3.2}

\begin{quote}\itshape
Perhaps a full model would require right-handed neutrinos. These would now be charged under \(U_A(1)_{L_{\mu-\tau}}\), which likely means the second and third generation must appear as vectorlike combinations. However, it seems to me that mixing between generations is prohibited by \(U_A(1)_{L_{\mu-\tau}}\).
\end{quote}

\noindent {\bf Reply.}

As the referee notes, the gauge charges also restrict the Dirac and Majorana mass terms when right-handed neutrinos are introduced. Adding such fermions while retaining the present symmetry-breaking sector does not remove the separation described above. In a type-I seesaw realization, the symmetry-enforced block structure is inherited by the effective light-neutrino mass matrix \(M_\nu=-m_DM_R^{-1}m_D^T\). A complete neutrino-sector extension would therefore require additional ingredients, such as a suitably charged scalar whose vacuum expectation value permits the missing mixing terms.

Neutrino-mass constructions involving right-handed neutrinos and additional symmetry-breaking scalars have been studied in the conventional gauged \(L_\mu-L_\tau\) framework; see Heeck and Rodejohann, Phys.\ Rev.\ D \textbf{84}, 075007 (2011) (\href{https://arxiv.org/abs/1107.5238}{arXiv:1107.5238}), and Asai, Hamaguchi and Nagata, Eur.\ Phys.\ J.\ C \textbf{77}, 763 (2017) (\href{https://arxiv.org/abs/1705.00419}{arXiv:1705.00419}). These provide relevant examples of neutrino-sector extensions, but are not identical to the axial construction considered here.

\medskip
\noindent {\bf Overall response to Major Point 3.}

\begin{quote}\itshape
Thus it is not obvious to me how neutrino oscillations can be included in this framework. If there is a simple solution, the authors should include it. If there is a known less-simple solution in the literature, the authors should cite this. Otherwise, at best this will be a confusing point. At worst, it would make the construction inconsistent with neutrino oscillation experiments.
\end{quote}

\noindent {\bf Reply.}

Developing such an extension would introduce additional fields and interactions beyond the model studied in this work. We therefore do not pursue a complete neutrino-mass and mixing mechanism here. Our statement is limited to the fact that the allowed effective operators can generate neutrino masses and \(\mu\)--\(\tau\) mixing; we do not claim that the present field content reproduces the full PMNS matrix. We thank the referee for drawing attention to this limitation and the need to distinguish these two issues.

To state this scope concisely in the manuscript, we propose the following clarification:
\begin{quote}
Neutrino masses and mixing within the \(\mu\)--\(\tau\) sector can be introduced through symmetry-allowed Weinberg-type operators, including higher-dimensional operators involving \(S\). However, the present field content does not generate the electron-flavour mixing required for the full observed PMNS matrix; a complete neutrino-mass and mixing mechanism would require additional ingredients and is not pursued in this work.
\end{quote}

\medskip
\noindent {\bf 4)}.
\begin{quote}\itshape
Is this the full set of experimental constraints? The second half of the manuscript presents a set of experimental bounds on the two-parameter effective \(A'\) model. It is not obvious if this is the full set of constraints on the model. (See also point \#1 above: I would have liked to see a section discussing the bounds on the \(S\) within this model since it cannot be decoupled.) The following experimental bounds seem like they could be relevant. It would be helpful to either include them or justify why they are subdominant to the experiments presented.
\end{quote}

\noindent {\bf Reply.}

We agree that the constraints presented for the two-parameter effective \(A'\) model do not, by themselves, establish a complete phenomenological assessment of the full UV construction. The relevant constraints must be separated by physical state and interaction structure. The scalar \(s\) constraints summarized in Table~II are not a substitute for stellar-cooling constraints on the gauge boson \(A'\); finite vector-like-lepton effects and kinetic mixing are also distinct from the strict pure-axial effective-current limit.

At present, the published scalar limits and their assumptions can be summarized, but the model-specific recasts of the scalar searches, the finite-mixing vector-versus-axial comparison, and the applicability of the cited \(A'\) stellar-cooling analysis have not all been established. We therefore do not claim that these constraints are globally subdominant or inapplicable. \textbf{[Insert the constraint-applicability summary and justified status for each constraint.]}

The scalar-sector constraints are addressed in connection with Major Point~1, while the four specific issues raised by the referee are discussed individually below.

\medskip
\noindent {\bf 4.1)}.
\begin{quote}\itshape
Neutrino oscillation observations. See point \#3 above. It appears that this model predicts a unit PMNS matrix, which is inconsistent with neutrino oscillations.
\end{quote}

\noindent {\bf Reply.}

This point is addressed in our response to Major Point~3. The allowed effective operators can generate neutrino masses and mixing within the \(\mu\)--\(\tau\) sector, so an identity PMNS matrix is not required by the symmetry. However, the electron-flavour mixing remains absent with the present field content, and the full observed PMNS matrix is not reproduced. A complete neutrino-mass and mixing mechanism would require additional ingredients and is not pursued in this work. The \(A'\) phenomenological results should be read under their stated interaction assumptions; they do not establish a complete explanation of the observed neutrino oscillations.

\medskip
\noindent {\bf 4.2)}.
\begin{quote}\itshape
There is no quantitative estimate of the amount of kinetic mixing generated with the ``benchmark'' choices of \(\theta\leq10^{-2}\). The authors make the insightful point that kinetic mixing goes to zero in the axial limit, but because this limit is tied to the size of the mixing (the mass of the vectorlike leptons), for the above benchmark there is non-zero kinetic mixing. It would be useful to have parametric estimates of the amount of kinetic mixing one expects so that readers may compare to experimental bounds. One could invoke parametric expressions in the literature for the relevant loop diagrams.
\end{quote}

\noindent {\bf Reply.}

We agree that a vanishing kinetic-mixing coefficient at a matching scale does not, by itself, quantify the residual mixing relevant to experiments. In the stated hypercharge--\(A'\) convention, the one-loop heavy-VLL threshold at zero momentum is
\[
\chi_{\rm EFT}(\mu_R)=\chi_{\rm UV}(\mu_R)
+\frac{g_Yg_XQ_\ell}{12\pi^2}
\ln\frac{M_{E_\mu}^2}{M_{E_\tau}^2}.
\]
It vanishes for the degenerate heavy-lepton benchmark \(M_{E_\mu}=M_{E_\tau}\), under the charge and coupling assumptions used in this matching result. A finite mass splitting produces a nonzero threshold even though the additive ultraviolet pole cancels.

This heavy-VLL hypercharge threshold must not be conflated with the low-energy photon--\(A'\) polarization induced by the \(\mu\)- and \(\tau\)-lepton loops, nor does the matching-scale condition \(\chi_{\rm UV}=0\) establish that every mixing contribution vanishes at experimental scales. The loop expression cited by the referee must be checked in the same charge, renormalization, and field-basis conventions before it is used for a numerical comparison.

\textbf{[Insert the verified loop estimate, matching scale and scheme, and comparison with applicable constraints.]}

\medskip
\noindent {\bf 4.3)}.
\begin{quote}\itshape
Another way of reframing the above point: the bounds presented are in the pure axial limit. Since the authors present a UV model that only realizes the axial coupling in a decoupling limit, it would be useful to comment on possible experimental bounds on the small vector-like coupling in a model with non-infinite vectorlike fermion masses. If these bounds are always small compared to the axial bounds, this should be stated and justified. The importance of this point is as follows: often bounds are sensitive to the chiral structure of interactions. Sometimes a vector interaction is more tightly constrained than its axial counterpart. Since the authors focus on the phenomenology in the axial limit, one wants to make sure that any small-but-finite vector interaction isn't significant because the constraints on the coupling are orders-of-magnitude larger than the axial coupling.
\end{quote}

\noindent {\bf Reply.}

We agree that the smallness of a residual vector coupling alone does not establish its phenomenological irrelevance, since experimental sensitivities can differ substantially between vector and axial interactions. In the conventions adopted in the manuscript, the muon couplings satisfy
\[
g_V^\mu=g_X\sin^2\theta_{\mu R},\qquad
g_A^\mu=-g_X\cos^2\theta_{\mu R},
\qquad
\left|\frac{g_V^\mu}{g_A^\mu}\right|
=\tan^2\theta_{\mu R}.
\]
Thus, for \(|\theta_{\mu R}|\leq10^{-2}\), the ratio is at most of order \(10^{-4}\). However, this parametric suppression alone is insufficient to establish that the residual vector interaction is subdominant to the axial interaction in the relevant experimental constraints.

The phenomenological implications of finite vector-like-lepton mixing must therefore be assessed by comparing the constraints applicable to the vector and axial interactions separately, while accounting for kinetic-mixing effects where relevant.

\textbf{[Insert the quantitative finite-mixing comparison and its impact on the plotted regions.]}

\medskip
\noindent {\bf 4.4)}.
\begin{quote}\itshape
Are stellar cooling bounds not relevant? See e.g. 2511.09619 for a recent paper on this. If they are not, then this is certainly worth justifying.
\end{quote}

\noindent {\bf Reply.}
We thank the referee for drawing our attention to this study. The current version of Blinov, Fox, Kelly, Plestid, and Zhou, \href{https://arxiv.org/abs/2511.09619}{arXiv:2511.09619} (v2, 15 May 2026), re-examines a sub-GeV gauge boson coupled to the \(L_\mu-L_\tau\) current. It discusses the standard SN~1987A cooling bound as well as diffusive cooling, constraints from low-energy supernovae, and the absence of a high-energy neutrino signal from SN~1987A. These results make clear that the supernova discussion should identify which physical probe and transport regime is being invoked, rather than treating “stellar cooling” as a single universal limit.

The scalar constraints summarized in Table~II do not, by themselves, determine whether the gauge boson \(A'\) is constrained in stellar environments. In addition, the cited analysis is for a vector \(L_\mu-L_\tau\) gauge boson, whereas the benchmark studied here has an axial charged-lepton current. Its exclusion curves therefore cannot be transferred directly without checking the current structure, neutrino and charged-lepton couplings, kinetic mixing and residual vector interactions, production rates, and transport assumptions. In particular, the relevant mass and coupling ranges and the conditions for trapping, escape, and decay must be compared explicitly.

A model-specific recast for the exact axial benchmark has not yet been established. We therefore do not claim that the cited stellar-cooling or supernova constraints are subdominant or inapplicable.

\textbf{[Insert the applicability assessment and recast status for arXiv:2511.09619.]}

\medskip
\noindent {\bf Minor Point 1)}. \begin{quote}\itshape
Range of parameters. The general construction of the model is clearly written, but it is difficult for the reader to envision the mass and coupling range that the authors are targeting. It would be very helpful if the authors decide on a range of \(m_{A'}\) and \(g_X\) of interest. For example: is there a range of masses where experimental or theoretical limits are less constraining? Is there some set of phenomenology that the authors seek to realize with a ``light'' \(A'\)? If so, they should define ``light.'' (See, e.g. the first word of the introduction, which is not defined with respect to any reference mass scale.)\end{quote}

\noindent {\bf Reply.}
We agree that the target parameter domain should be specified and physically motivated, and that the adjective ``light'' should not be used without a reference scale. The intended revision is to replace the opening sentence of the Introduction with ``Neutral spin-one bosons'' rather than the unqualified phrase ``Light neutral spin-one bosons''. \textbf{[Insert the final Introduction section reference after applying the edit.]}

The parameter scan shown in Fig.~2 covers
\[
10^{-3}\,\mathrm{GeV}\leq m_{A'}\leq10^2\,\mathrm{GeV},
\qquad
10^{-6}\leq g_X\leq10^{-1}.
\]
These intervals define the plotting window used in the phenomenological analysis; they do not, by themselves, represent theoretical boundaries or experimentally allowed regions. \textbf{[Insert the physical rationale and references for the plotted scan ranges.]}

\medskip
\noindent {\bf Minor Point 2)}.
\begin{quote}\itshape
Defining the pure-axial limit. Equation (12) gives the coupling of the \(A'\) to the mass-basis fields. It is here that the ``pure axial coupling'' limit is defined? i.e. \(\theta_{\mu R}\to0\)? If that is the case, it would help to define this as the pure-axial limit, a phrase which is used below (36) but not otherwise quantified.
\end{quote}

\noindent {\bf Reply.}

We agree that the pure-axial limit should be defined explicitly when the mass-basis couplings are introduced. We define the strict pure-axial limit as \(\theta_{\mu R}\to0\), for which the muon couplings
\[
g_V^\mu=g_X\sin^2\theta_{\mu R},
\qquad
g_A^\mu=-g_X\cos^2\theta_{\mu R}
\]
reduce to \(g_V^\mu=0\) and \(g_A^\mu=-g_X\), respectively. Thus, the coupling of \(A'\) to the muon becomes purely axial in this limit. The same limiting convention applies to the \(\tau\) sector with the charge reversal specified in the model.

For finite \(\theta_{\mu R}\), the interaction contains a nonzero vector component and is therefore not strictly purely axial. The term ``pure-axial limit'' is reserved for this limiting case; finite \(\theta_{\mu R}\) gives a nonzero vector component. \textbf{[Insert the final section and equation references.]}

\medskip
\noindent {\bf Minor Point 3)}.
\begin{quote}\itshape
Bounds on the vectorlike leptons. Because the mixing angles scale inversely with the vectorlike lepton masses, it would be helpful to present lower limits on these masses. This, in turn, gives an upper limit on the \(\theta_{(\mu,\tau)R}\) mixing angles which then gives an upper limit on the \(A'\) vectorlike coupling to \(L_{\mu-\tau}\). How do these bounds compare with the \(\theta_{(\mu,\tau)(L,R)}\) bounds in equation (27)?
\end{quote}

\noindent {\bf Reply.}

We agree that the vector-like-lepton mass constraints, the bounds on the light--heavy mixing angles, and the resulting \(A'\) couplings should be presented consistently, while distinguishing theoretical relations from direct experimental limits. In the seesaw-like regime, the charged-lepton mass relations in Eqs.~(21) and (23) give
\[
m_\ell\simeq\theta_{\ell L}\theta_{\ell R}M_{E\ell},
\qquad \ell=\mu,\tau.
\]
Combining these relations with the adopted small-mixing condition in Eq.~(27),
\[
|\theta_{\ell L}|,\;|\theta_{\ell R}|\leq10^{-2},
\]
yields the following conditional mass scales at leading order:
\[
M_{E\mu}\gtrsim
\frac{m_\mu}{(10^{-2})^2}
\simeq1.06~\mathrm{TeV},
\qquad
M_{E\tau}\gtrsim
\frac{m_\tau}{(10^{-2})^2}
\simeq17.8~\mathrm{TeV}.
\]
For the degenerate-mass benchmark \(M_{E\mu}=M_{E\tau}\equiv M_E\) adopted in Eq.~(22), the tau-sector relation gives the stronger conditional requirement \(M_E\gtrsim17.8~\mathrm{TeV}\). This is a consequence of combining the leading-order mass relations with the adopted mixing-angle condition, rather than a direct experimental lower bound. This conditional mass scale is distinct from a direct collider lower bound. The condition \(M_E>10~\mathrm{TeV}\) in Eq.~(22), considered by itself, is weaker than the leading-order conditional requirement above.

The mass relations constrain the products \(\theta_{\ell L}\theta_{\ell R}\), rather than the two mixing angles independently. Consequently, a direct-search lower bound on \(M_{E\ell}\) constrains the corresponding mixing-angle product; an individual upper bound on \(\theta_{\ell R}\) requires additional information on \(\theta_{\ell L}\) or an explicit assumption about the relative sizes of the two mixing angles. The individual mixing-angle constraints in Eq.~(27) must therefore be distinguished from mass limits inferred from the charged-lepton mass relations.

The implications for the \(A'\) interactions follow from the mass-basis couplings. In the conventions adopted in the manuscript,
\[
g_V^\mu=g_X\sin^2\theta_{\mu R},
\qquad
g_A^\mu=-g_X\cos^2\theta_{\mu R},
\qquad
\left|\frac{g_V^\mu}{g_A^\mu}\right|
=\tan^2\theta_{\mu R}.
\]
Thus, the adopted condition \(|\theta_{\mu R}|\leq10^{-2}\) implies a relative vector component of at most approximately \(10^{-4}\). This is a consequence of the mixing-angle assumption, not an independent experimental bound, and does not by itself establish the phenomenological irrelevance of the residual vector interaction.

\textbf{[Insert the validated VLL mass and mixing limits, assumptions, and references.]}

\medskip
\noindent {\bf Minor Point 4)}.
\begin{quote}\itshape
Conservative benchmark. Equation (27) is presented as a ``conservative benchmark.'' What does this mean? A benchmark point would be a choice of these mixing angles. A constraint is an upper bound. Eq. (27) is written as an upper bound, but is described as a benchmark. This is not simply semantic: readers would want to understand if \(10^{-2}\) is an experimental upper limit or simply a convenient choice.
\end{quote}

\noindent {\bf Reply.}
We agree that the term ``conservative benchmark'' was ambiguous because Eq.~(27) specifies upper bounds on the mixing angles rather than a particular benchmark point. The appropriate terminology is ``small-mixing assumption'', because
\[
|\theta_{\ell L}|,\;|\theta_{\ell R}|\leq10^{-2},
\qquad \ell=\mu,\tau,
\]
is an assumption adopted for the phenomenological analysis, rather than a universal experimental upper limit on the mixing angles. The purpose of this assumption is to specify the small-mixing regime considered in our analysis; it does not, by itself, establish that the chosen values are conservative with respect to all relevant experimental constraints.

\textbf{[Insert the applicable mixing-angle constraints and the comparison with the adopted criterion.]}

\medskip
\noindent {\bf Minor Point 5)}.
\begin{quote}\itshape
What is the purpose of including muon collider bounds? The second half of the paper focuses on existing experimental bounds on the effective theory. It is a strange change in tone to include hypothetical muon collider bounds for future reach. The authors do not justify why a future muon collider is the way to test this model in the future and the section feels like it is tacked on. Most of the section emphasizes that the muon collider search would be a ``standard'' follow up of existing work (Ref 83) and so this section does not seem particularly novel. The authors simply focus on the forward-backward asymmetry. Why not also include polarization asymmetries? Or would there be asymmetries in the \(ee\) versus \(\tau\tau\) final states? Given the first half of the paper, though, it seems that the real ``smoking gun'' for this UV completion is the physics of the \(S\) particle which should be tested in any Higgs-factory phase of a future lepton collider. I feel that this section either adds too little and should be downplayed as an example of future tests, or otherwise should have its own full fledged exploration (which may be a separate paper).
\end{quote}

\noindent {\bf Reply.}

We thank the referee for this constructive comment. We agree that the scope of the future-muon-collider discussion should match the analysis actually performed. The proposed revision treats the section as an illustrative outlook on possible future tests of the effective \(A'\) interactions, not as a comprehensive collider forecast or a projected exclusion/discovery reach.

The forward--backward asymmetry is retained only as an illustrative chirality-sensitive observable within the assumptions of the present analysis. The paper does not claim to have quantitatively studied polarization asymmetries or comparisons among \(ee\) and \(\tau\tau\) final states. A future lepton collider's Higgs-factory program may also probe the scalar sector, but without a dedicated model-specific sensitivity analysis no quantitative claim is made here about \(S\) production rates, decay signatures, or reach.

\textbf{[Insert the agreed Muon Collider scope and final section reference.]}

\medskip
\noindent {\bf Minor Point 6)}. \begin{quote}\itshape
Terminology. The following points could be improved for clarity.
\end{quote}

\medskip
\noindent {\bf Minor Point 6.1)}. \begin{quote}\itshape
The word ``light'' on page 2 has no reference scale.
\end{quote}

\noindent {\bf Reply.}
We agree that the term ``light'' should not be used without an explicit reference scale. We have replaced the opening phrase ``Light neutral spin-one bosons'' with ``Neutral spin-one bosons'' in Sec.~I. The Abstract continues to state the explicit MeV-to-TeV mass range.

\medskip
\noindent {\bf Minor Point 6.2)}. \begin{quote}\itshape
At the bottom of page 11, when referring to Table II, the authors write ``representative bounds.'' What does representative mean? Does it mean these are the most stringent bounds? Are these ``some'' bounds? Is the list ``complete'' in any sense?
\end{quote}

\noindent {\bf Reply.}
We agree that ``representative bounds'' could be interpreted as implying either a complete compilation or the most stringent constraint at every scalar mass. The table reproduced under Major Point~1 summarizes selected published constraints on the scalar--muon coupling \(y_{s\mu\mu}\), including neutron-star dynamics, fifth-force searches, neutron-star cooling, and SN~1987A cooling and diffuse gamma rays. It is not an exhaustive inventory and does not imply that the listed bounds are uniformly the strongest at every mass.

The quoted limits inherit the production, trapping, and loop-matching assumptions of the corresponding analyses. The revised manuscript should describe them as selected, non-exhaustive published constraints, not as a complete or model-independent exclusion of the scalar parameter space. \textbf{[Insert the final Table II caption and text references after applying the edit.]}

\medskip
\noindent {\bf Minor Point 6.3)}. \begin{quote}\itshape
Between page 12--13: ``...do not provide a universal \(g_X\) bound, but they show that the intermediate scalar-mass region cannot simply be assumed unconstrained.'' What does intermediate mass mean? It is also not clear what the intended content of this statement is: are the authors presenting a bound? Why would the reader assume that a new particle in this range would be unconstrained?
\end{quote}

\noindent {\bf Reply.}
We agree that the phrase ``intermediate scalar-mass region'' was insufficiently precise and did not clearly distinguish experimental search coverage from a model-specific exclusion. The relevant mass ranges are those probed by the published searches for leptophilic scalars, rather than a separately defined universal mass interval.

BaBar and Belle searched for
\(e^+e^-\to\tau^+\tau^-s\), followed by
\(s\to e^+e^-\) or \(\mu^+\mu^-\), over approximately
\(0.04\,\mathrm{GeV}\lesssim m_s\lesssim6.5\)--\(7\,\mathrm{GeV}\).
The Belle~II four-muon resonance search (arXiv:2403.02841) includes a muonphilic spin-0 interpretation over approximately \(0.212\)--\(9\,\mathrm{GeV}\); the \(\mu^+\mu^-\tau^+\tau^-\) resonance search (arXiv:2306.12294) includes a Higgs-like leptophilic spin-0 interpretation over approximately \(3.6\)--\(10\,\mathrm{GeV}\) in the \(\tau^+\tau^-\) resonance mass. These are distinct analyses with different production channels, acceptances, and scalar benchmarks; neither mass coverage alone nor the existence of a spin-0 interpretation establishes a constraint on our UV scalar without checking the model-specific couplings and branching fractions.

For the mass-proportional scalar benchmark used in the Belle analysis, the published limit \(\xi_{\rm lim}=0.483\) at \(m_s=1\,\mathrm{GeV}\) implies, at 90\% confidence level,
\[
v_s>\frac{v_h}{\xi_{\rm lim}}
\simeq509\,\mathrm{GeV},
\qquad
y_{s\mu\mu}<2.1\times10^{-4},
\]
where \(v_h=246\,\mathrm{GeV}\). Using \(m_{A'}=2g_Xv_s\), this gives the conditional relation
\[
g_X<9.8\times10^{-4}
\left(\frac{m_{A'}}{\mathrm{GeV}}\right).
\]
For this scalar-mass benchmark, the direct application of the visible-decay limit requires the \(s\to A'A'\) channel to be kinematically closed, or its effect on the visible branching fractions to be accounted for. The resulting relation is therefore a benchmark-dependent constraint, not a universal exclusion for the full UV completion.

In our model, the decay \(s\to A'A'\) becomes kinematically allowed when \(m_s>2m_{A'}\). When this channel is open, it can modify the visible scalar branching fractions, so the published visible-decay limits cannot be transferred directly without accounting for this additional decay mode.

The revised text should replace the phrase ``intermediate scalar-mass region'' with explicit mass ranges and a distinction between published search coverage and a model-specific exclusion. The recast status for the present UV scalar remains \textbf{[Insert the model-specific scalar-recast status after checking production and decay assumptions.]}; no region is classified as excluded or allowed solely from mass coverage.

\medskip
\noindent {\bf Minor Point 6.4)}. \begin{quote}\itshape
Below (34): ``we impose Eq. (34) point by point'', it's not clear what is meant by ``point.''
\end{quote}

\noindent {\bf Reply.}
A ``point'' denotes one grid point \((m_{A'},g_X)\) in the parameter scan shown in Fig.~2. At each point, the symmetry-breaking scale is derived from
\[
v_s=\frac{m_{A'}}{2g_X},
\]
rather than varied independently. For a specified scalar mass \(m_s\), the scalar self-coupling is then determined by
\[
\lambda_S
=\frac{m_s^2}{2v_s^2}
=\frac{2g_X^2m_s^2}{m_{A'}^2}.
\]
The condition
\[
\lambda_S\leq\lambda_S^{\max}
\]
is consequently checked at each scan point for the adopted scalar benchmark.

This point-by-point check tests the theoretical consistency of the chosen scalar benchmark. It is distinct from the application of experimental constraints, which must be assessed separately. The manuscript should state this distinction next to Eq.~(34).

\medskip
\noindent {\bf Minor Point 6.5)}. \begin{quote}\itshape
Page 21, first paragraph of Section E. What is the ``parameter region of interest''? This wasn't clearly defined anywhere.
\end{quote}

\noindent {\bf Reply.}
We agree that the phrase ``parameter region of interest'' did not specify the mass coverage of the searches considered in Section~III.E. The relevant four-muon searches cover approximately
\[
\begin{aligned}
\text{BaBar}:&\quad 0.212\text{--}10\,\mathrm{GeV},\\
\text{Belle II}:&\quad 0.212\text{--}9\,\mathrm{GeV},\\
\text{CMS}:&\quad 5\text{--}70\,\mathrm{GeV}.
\end{aligned}
\]
These intervals describe the mass coverage of the published analyses; they do not imply that the entire intervals are excluded in our model.

The relevant processes are
\[
e^+e^-\to\mu^+\mu^-A',
\qquad A'\to\mu^+\mu^-,
\]
and
\[
Z\to\mu^+\mu^-A',
\qquad A'\to\mu^+\mu^-.
\]
The reinterpretation accounts for the production rate and \(\mathrm{Br}(A'\to\mu^+\mu^-)\). The experimental analyses also impose prompt-resonance selection requirements. Their limits are applicable to our model only where the predicted \(A'\) decay length is compatible with those requirements; the mass coverage alone does not establish this condition throughout the parameter space.

As an illustrative check, at \(m_{A'}=0.212\,\mathrm{GeV}\) and \(g_X=7\times10^{-4}\), the neutrino decay channels alone give \(c\tau_{A'}\lesssim0.072\,\mu\mathrm{m}\), corresponding to a laboratory-frame decay length of approximately \(1.8\,\mu\mathrm{m}\) for \(\gamma\simeq25\). This estimate supports the prompt-decay approximation at this benchmark point, but does not establish its validity throughout the full parameter space.

The vector-to-axial rate-matching approximation is controlled for \(m_{A'}\gg m_\mu\). Near the dimuon threshold, the rate-level comparison does not by itself establish a detector-level exclusion. The revised text should state the experimental mass coverage and the assumptions underlying the reinterpretation, without extrapolating prompt-resonance constraints beyond the regions where their applicability is established. The displayed decay-length estimate is a single benchmark check, not a validation of the full scan.

\medskip
\noindent {\bf Minor Point 6.6)}. \begin{quote}\itshape
Page 21, end of section D: What's the significance of the \(v_s=2\,\mathrm{GeV}\) line? What does it mean that \(y_{s\mu\mu}\) ``becomes large'' (nonperturbative)? Why plot it? Is it a perturbativity bound?
\end{quote}

\noindent {\bf Reply.}
The line at \(v_s=2\,\mathrm{GeV}\) is a theoretical reference associated with the benchmark \(m_s=10\,\mathrm{GeV}\) and \(\lambda_S^{\max}=4\pi\). From
\[
m_s^2=2\lambda_Sv_s^2,
\]
the scalar self-coupling condition gives
\[
v_s\geq\frac{m_s}{\sqrt{2\lambda_S^{\max}}}
=\frac{10\,\mathrm{GeV}}{\sqrt{8\pi}}
\simeq1.995\,\mathrm{GeV}.
\]
Thus, for this benchmark, the line marks the approximate lower boundary on \(v_s\); values below this boundary violate the adopted scalar self-coupling criterion. The line is not an experimental exclusion and is not a perturbativity bound on \(y_{s\mu\mu}\) itself.

Using the coupling relation adopted in our manuscript,
\[
y_{s\mu\mu}=\frac{m_\mu}{v_s},
\]
we obtain
\[
y_{s\mu\mu}\simeq
\frac{0.10566\,\mathrm{GeV}}{2\,\mathrm{GeV}}
\simeq0.053
\qquad\text{at}\qquad v_s=2\,\mathrm{GeV}.
\]
A useful measure of the perturbative expansion is
\[
\frac{y_{s\mu\mu}^2}{4\pi}
\simeq2.2\times10^{-4}\ll1.
\]
Thus, the Yukawa coupling itself is well within the perturbative regime at this reference value of \(v_s\). The phrase ``becomes large'' should therefore be understood as describing the increase of \(y_{s\mu\mu}\) as \(v_s\) decreases, rather than indicating that the coupling becomes nonperturbative at \(v_s=2\,\mathrm{GeV}\).

The line at \(v_s\simeq2\,\mathrm{GeV}\) instead corresponds to the adopted scalar self-coupling criterion for \(m_s=10\,\mathrm{GeV}\) and \(\lambda_S^{\max}=4\pi\). The discussion and figure caption should identify this line as a benchmark-dependent scalar-sector consistency boundary, not as a perturbativity bound on \(y_{s\mu\mu}\) or an experimental exclusion contour. \textbf{[Insert the final manuscript reference after applying the edit.]}

\medskip
\noindent {\large Summary:}

We thank the referee again for the careful and constructive report. The revised manuscript must state the assumptions underlying the effective-current phenomenology, distinguish the plotted parameter window from theoretical and experimental boundaries, and separate established constraints from benchmark-dependent interpretations and regions that remain unassessed.

Regarding neutrino oscillations, the symmetry-allowed effective operators can generate neutrino masses and mixing within the \(\mu\)--\(\tau\) sector, but the present field content does not generate the electron-flavour mixing required for the full observed PMNS matrix. Our replies to Major Points~3 and 4.1 distinguish these two statements. We do not pursue a complete neutrino-mass and mixing mechanism here; this would require additional fields and interactions beyond the model studied in this work. A concise manuscript clarification is proposed in the overall response to Major Point~3.

Regarding the scalar \(S\), we retain it as a physical state in the UV construction and do not claim that its mixing with the Higgs is negligible over the full scalar-mass range. The zero-momentum portal threshold and the explicit small-portal matching condition are stated as benchmark inputs. The response now reports the completed benchmark-level finite-momentum one-loop fermion self-energies and the width-regulated near-degeneracy diagnostic, while distinguishing these from a physical mixing angle or a global viability result. The full pole-scheme treatment, diagonal self-energies and tadpole/input counterterms, and model-specific two-pole Higgs lineshape and direct-search recasts remain open; the narrow-width near-degenerate region therefore remains \textbf{not assessed}.

The effective-current limits are interpreted only under the interaction structures and assumptions for which their applicability has been established. The final constraint classification is \textbf{[Insert the final kinetic-mixing, finite-vector, and stellar-cooling applicability classification.]}

The referee concludes:

\begin{quote}
\itshape
In its current form, the paper is not ready for publication. The major points above must be addressed. It's entirely possible that I have misinterpreted some of these points, but if that is the case then they at least indicate ideas that must be clarified before the manuscript merits publication. I strongly encourage the authors to further address the minor points.

The paper can be a useful addition to the model building literature. However, the question of whether this class of model can be consistent with neutrino oscillations is critical. I also feel that dismissing the physics of the scalar \(S\) ignores the actual predictions of the UV model---otherwise the authors are simply justifying the phenomenology of the effective axial theory.
\end{quote}

We thank the referee for emphasizing that the scalar \(S\) is a physical prediction of the UV construction. We do not discard \(S\): we state its symmetry-breaking role and mass relations, summarize the relevant constraints, and cite the extensive literature on singlet--Higgs mixing and scalar searches. Our conclusions on the \(A'\) phenomenology are restricted to the explicitly stated small-portal benchmark; the conditional light-lepton estimate is not presented as proof that mixing is negligible over the full scalar-mass range. \textbf{[Insert the final section references after applying the manuscript revisions.]}

\medskip

\noindent Sincerely yours,

\noindent Rundong Fang, Jinhui Guo, Ming Li, Jia Liu, Xiao-Ping Wang, Yiheng Xiong

\bibliographystyle{JHEP}
\bibliography{references}
\end{document}